% ============================================================
% Chapter 1: Introduction
%   Section: Context
% Included from chapters/01-introduction/introduction.tex
% Source: docs/resumen-ejecutivo.md, "El problema que resuelve";
% docs/idea-de-negocio.md, section 1.
% ============================================================

\section{Context}
\label{sec:context}

Three problems of today's news meet in this work.

\paragraph{False information travels faster than its correction.}
On Twitter between 2006 and 2017, false stories reached more people,
faster and further than true ones, and the difference was driven by
people rather than by automated accounts \cite{vosoughi2018spread}.
Professional fact-checking grew in response: by the middle of the last
decade dozens of fact-checking organisations were active in Europe
alone \cite{graves2016rise}. Their method is manual --- a trained
person chooses a claim, searches for evidence and writes a verdict
--- and it cannot keep pace with the volume of what is published.
Automated fact-checking, the research field that tries to support or
partly automate that work, frames it as a pipeline of claim detection,
evidence retrieval and verdict prediction \cite{guo2022survey}, and
increasingly as a tool to assist fact-checkers rather than to replace
them \cite{nakov2021assisting}.

\paragraph{Distribution optimises for engagement.} The systems that
decide what a reader sees are usually optimised for clicks and
attention, not for accuracy or usefulness, and personalisation tuned
that way tends to narrow what each reader is shown
\cite{pariser2011filter,karimi2018news}.

\paragraph{The news cycle is dominated by negative content.} Readers
react more strongly to negative news than to positive news
\cite{soroka2015negativity}, and a news supply shaped by that reaction
erodes trust without necessarily giving readers anything they can act
on. \emph{Constructive journalism} answers by reporting problems
together with the responses to them, keeping journalism's core
functions while applying what positive psychology has learned about
how news affects its audience \cite{mcintyre2017constructive}.

The project behind this thesis starts from the combination of the
three: a source of constructive news on chosen topics whose claims are
checked against real evidence before they are published, and whose
verdicts can be inspected by whoever reads them. This thesis builds
and evaluates the first part of that project --- finding the news and
verifying it --- and leaves its distribution to future work.
